En un ensayo simulado, una concentración \(c(t)>0\), en \(\mathrm{g/L}\), satisface \(c'=-kc\), con \(k>0\) y \(t\) en horas. No se conoce \(c(0)\). El instrumento de calibración informa de la media temporal
\[
\overline c_{[u,v]}=\frac1{v-u}\int_u^v c(t)\,dt,
\]
no de una concentración instantánea. Sus dos lecturas, tratadas como exactas, son
\[
\overline c_{[0,1]}=0.800\,\mathrm{g/L},\qquad
\overline c_{[1,2]}=0.400\,\mathrm{g/L}.
\]
Se ha reservado una medición instantánea en \(t=3\) h: \(c_{\rm obs}=0.140\,\mathrm{g/L}\), con error absoluto garantizado no mayor que \(0.004\,\mathrm{g/L}\).

Infiera \(k\) y \(c(0)\) a partir de las dos medias y prediga la medición reservada. Compruebe su compatibilidad con la cota de error. Explique por qué no sería válido sustituir la media de cada intervalo por la concentración en su extremo derecho al inferir el estado inicial.
